\section*{Erd\H{o}s Problem 1202: clustered primes defeat the proposed sieve}
\subsection*{1. Statement and outcome}
Does each $\epsilon>0$ admit a $k$ such that, for every choice of $k$ primes below $n^{1-\epsilon}$ and $(p-1)/2$ forbidden residue classes at each prime $p$, at most $\epsilon n$ integers in $[1,n]$ survive? The printed statement also introduces $\eta>0$ but gives it no subsequent role. For that recorded statement the answer is negative. We give a construction with $\epsilon=1/4$, arbitrarily many primes of order $n^{2/3}$, and more than $0.48n$ survivors for sufficiently large $n$.
The known negative resolution is attributed to Price and GPT-5.4 Pro. The following clustered-prime construction spells out a sufficient counterexample; no novelty is asserted.

\subsection*{2. Finding many primes in a narrow interval}
We use the prime number theorem in the elementary consequence that, for all sufficiently large $m$, the number of primes in $[m,2m)$ is at least $c m/\log m$ for some absolute $c>0$. No prime number theorem in short intervals is assumed.

Given a large integer $n$, put
\[
 m=\lfloor n^{2/3}\rfloor,\qquad H=\left\lfloor\frac{m^2}{100n}\right\rfloor.
\]
Then $H\to\infty$. Partition the integer interval $[m,2m)$ into at most $\lceil m/H\rceil$ bins of length at most $H$. By averaging, one bin, with integer left endpoint $Q\in[m,2m)$, contains
\[
 \gg\frac{H}{\log m}\gg\frac{n^{1/3}}{\log n}
\]
primes, all in $[Q,Q+H)\cap[m,2m)$. In particular the number available tends to infinity. For any prescribed $k$, choose $n$ large enough and select $k$ distinct primes from this bin. Since $2m<n^{3/4}$ for sufficiently large $n$, they satisfy the required bound $p<n^{1-\epsilon}$ with $\epsilon=1/4$.

\subsection*{3. A common surviving set}
Set
\[
 B=\left\lceil\frac{nH}{m}\right\rceil,\qquad
 T=\left\lfloor\frac nQ\right\rfloor,
\]
and consider
\[
 S=\left\{aQ+b:0\le a<T,\quad B\le b\le
 \left\lfloor\frac{Q-1}{2}\right\rfloor\right\}. \tag{1}
\]
For large $n$ the interval of allowed $b$ is nonempty and $B\ge1$. All members lie in $[1,n]$, and distinct pairs $(a,b)$ give distinct integers, because $0\le b<Q$.

For each chosen odd prime $p$, declare
\[
 A_p=\{(p+1)/2,(p+3)/2,\ldots,p-1\}
\]
to be the forbidden residues. This has exactly $(p-1)/2$ members. Write $p=Q+\delta$, where $0\le\delta<H$. If $z=aQ+b\in S$, then
\[
 z=ap+(b-a\delta),\qquad
 0\le b-a\delta\le b\le\left\lfloor\frac{Q-1}{2}\right\rfloor
 \le\frac{p-1}{2}. \tag{2}
\]
Indeed $a<n/Q\le n/m$ and $b\ge B\ge nH/m\ge a\delta$. Thus the parenthesized quantity is the actual least nonnegative residue modulo $p$, and lies in the allowed lower half. Equation (2) holds for every selected prime, so every member of $S$ survives all $k$ exclusions simultaneously.

\subsection*{4. Counting and quantifiers}
Since $H\le m^2/(100n)$ and $Q\ge m$,
\[
 \frac BQ\le\frac1{100}+\frac1Q.
\]
The number of choices of $b$ in (1) is at least $Q/2-B$. Consequently, for sufficiently large $n$,
\[
 \begin{split}
 |S|&\ge\left(\frac nQ-1\right)(Q/2-B)\\
 &\ge \frac n2-\frac{nB}{Q}-\frac Q2
 \ge0.49n-\frac nQ-\frac Q2
 =0.49n-O(n^{2/3}).
 \end{split}
\]
In particular $|S|>0.48n>n/4$ eventually. For every fixed $k$ there are arbitrarily large $n$ for which the bin supplies $k$ primes and all these inequalities hold. Hence no proposed $k$ works for $\epsilon=1/4$. Any value of the unused parameter $\eta$ can be fixed without affecting the example.

The full negative answer to the printed statement is proved relative to the stated prime number theorem input. The construction places the primes above the square-root scale, so it is compatible with the positive large-sieve observation for primes below $\sqrt n$.
Source and attribution: \url{https://www.erdosproblems.com/1202}, including its account of the Price--GPT-5.4 Pro resolution and the associated discussion at \url{https://www.erdosproblems.com/forum/thread/1202}.

\textbf{Completion rate estimate: 100\%.}
